% !TEX root = ../main.tex
\section{Related Work}
\label{sec:background}

This section situates our study within two bodies of prior work: empirical
studies of software marketplaces and platform governance, and research on
interaction modality and dual-audience design.

%% ─────────────────────────────────────────────
\subsection{Software Marketplaces and Platform Governance}
\label{sec:rw:marketplace}

A substantial body of work examines software marketplaces, primarily mobile app
stores: sales rank strongly shapes mobile-app demand
\cite{ghose2014apppricing}, platform-owner entry into complementor markets
induces complementor responses
\cite{wen2019platformowner}, and app updates and category positioning predict
adoption better than review scores \cite{janssen2022googleplay}, while update
management is a strategic lever that differs across app stores
\cite{comino2019appstore}. More recently,
\citet{su2024gptstore} documents the OpenAI GPT Store's extreme concentration
and ``supply without demand,'' and \citet{stein2026mcpecosystem} documents the
MCP ecosystem's reshaping of developer effort. Platform governance research
\cite{rietveld2021platform} explains how such outcomes arise: certification can
amplify concentration by disproportionately benefiting already-popular apps
\cite{rietveld2021marketorchestrators};
boundary resources mediate the owner--complementor relationship
\cite{eaton2015governance,tiwana2015platformecosystems}; and two-sided market
theory \cite{rochet2003twosided} and network effects \cite{caillaud2003chicken}
explain winner-take-all dynamics, with \citet{boudreau2012complementors} showing
that more complementors can reduce individual incentives even as aggregate
value rises. Concentration is standardly measured with Gini coefficients and
HHI, as in analyses of the GPT Store \cite{su2024gptstore}, and the
``superstar'' phenomenon is well-documented
\cite{aguiar2018platforms}. These studies establish heavy-tailed adoption and
platform design effects, but none examines automation marketplaces, or whether
design decisions that introduce new consumption modalities (such as agentic
access) are associated with redistributed concentration.

%% ─────────────────────────────────────────────
\subsection{Agents as Users and Dual-Audience Design}
\label{sec:rw:agentic}

The term ``agentic economy'' describes the emerging ecosystem in which AI
agents act as autonomous consumers of digital services---discovering tools,
negotiating access, executing tasks, and making payments without continuous
human oversight. Governing such agentic systems is an emerging policy concern
\cite{shavit2023practices}. This shift is enabled by large
language models (LLMs) with tool-use capabilities
\cite{schick2023toolformer,qin2024toolllm}, standardized tool-calling protocols
such as MCP---now spanning a large tool ecosystem
\cite{stein2026mcpecosystem}---and JSON-schema-based function calling
\cite{openai2024function}, and payment infrastructure that permits
machine-initiated transactions.

For automation marketplaces, the agentic transition is a change in
\emph{interaction modality}: where a human once browsed a catalog, filled a
form, and clicked run, an agent now discovers via semantic search, generates
parameters from a schema, and pays through pre-authorized billing. These
changes may favor tools that are well-documented and schema-rich.

We treat this modality transition as the paper's theoretical object of study,
rather than as background context, and ground it in four intersecting HCI
lineages.

\emph{Human--AI interaction.} Work on human--AI collaboration establishes
guidelines for systems that collaborate with users \cite{amershi2019guidelines}
and examines how humans and agents negotiate roles beyond a linear
``user specifies, agent executes'' model \cite{zhou2024nonlinear}. Our setting
asks what changes when the ``user'' is itself an agent.

\emph{Agents as users, and machine legibility.} A growing literature treats AI
agents as end-users with their own needs: interface design alone drives large
performance gains for agents \cite{yang2024sweagent}. API usability research
establishes that interfaces demand explicit design for their users
\cite{myers2016api}; extending this principle to machine consumers suggests
that interfaces built for one audience can become illegible to another. These
motivate our question of whether a \emph{payment-eligibility signal} is
associated with measurable adoption differences.

\emph{Developer experience and mediated work.} Developer experience (DevEx) work
\cite{fagerholm2012developer,weisz2025codeassistant} and algorithmic-mediation
research \cite{eslami2015algorithmic,lee2015working} show that platforms
designed for one population can be illegible to another, and recommendation
systems have been framed as requiring fairness across multiple stakeholder
groups \cite{burke2017multisided}---direct concerns for a marketplace whose
audience is now heterogeneous.

\emph{Hybrid intelligence and automation paradigms.} Hybrid-intelligence
research frames human--AI collaboration as a division of labor
\cite{dellermann2019hybrid}, and automation-paradigm studies show that
``do it for me'' versus ``do it with me'' is a design decision with distinct
consequences for control and trust \cite{khurana2025doitforme}. Our paper asks
what a design decision optimizing for one modality does to the other.

Apify's agentic-readiness boundary is a concrete instance of such a design
decision: it admits tools with limited permissions and no standby mode to
agent-initiated use. Whether this readiness is associated with measurable
adoption differences is the empirical question this paper studies.
